% answers.tex —— 2005 年·普通高等学校招生考试·重庆卷·理科 · 参考答案与详细解析片段
% 由 tools/generate_projects.py 依来源总账自 2005/2005_数学真题_重庆理科_试题.md + 2005/2005_数学真题_重庆理科_解析正文.md + 2005/2005_数学真题_重庆理科_答案要点.md 生成，勿手改（改源 MD 后重跑）

\begin{answer}
\gksolhead{1}{}

\textbf{题目}：

圆 $(x + 2)^2 + y^2 = 5$ 关于原点 $(0, 0)$ 对称的圆的方程为\mbox{（　　）}

\begin{gkopts}(2)
\task A．$(x - 2)^2 + y^2 = 5$
\task B．$x^2 + (y - 2)^2 = 5$
\task C．$(x + 2)^2 + (y + 2)^2 = 5$
\task D．$x^2 + (y + 2)^2 = 5$
\end{gkopts}

\Needspace{4\baselineskip}
\textbf{解析}

原圆方程为 $(x + 2)^2 + y^2 = 5$，其圆心坐标为 $(-2, 0)$，半径为 $\sqrt{5}$．
圆心 $(-2, 0)$ 关于原点 $(0, 0)$ 的对称点坐标为 $(2, 0)$．
因为对称后的圆半径保持不变，仍为 $\sqrt{5}$，
所以所求圆的方程为 $(x - 2)^2 + y^2 = 5$．
故选 A．

\gkanswer{答案：A}

\end{answer}
\begin{answer}
\gksolhead{2}{}

\textbf{题目}：

$\left(\dfrac{1 + \mathrm{i}}{1 - \mathrm{i}}\right)^{2005} =$\mbox{（　　）}

\begin{gkopts}(4)
\task A．$\mathrm{i}$
\task B．$-\mathrm{i}$
\task C．$2^{2005}$
\task D．$-2^{2005}$
\end{gkopts}

\Needspace{4\baselineskip}
\textbf{解析}

\Needspace{5\baselineskip}
首先化简括号内的复数分式：
\nopagebreak[4]
\[
\frac{1 + \mathrm{i}}{1 - \mathrm{i}} = \frac{(1 + \mathrm{i})^2}{(1 - \mathrm{i})(1 + \mathrm{i})} = \frac{1 + 2\mathrm{i} - 1}{1 - (-1)} = \frac{2\mathrm{i}}{2} = \mathrm{i}.
\]
因为虚数单位 $\mathrm{i}$ 的乘幂具有周期性 $\mathrm{i}^4 = 1$，且 $2005 = 4 \times 501 + 1$，
\Needspace{5\baselineskip}
所以：
\nopagebreak[4]
\[
\left(\frac{1 + \mathrm{i}}{1 - \mathrm{i}}\right)^{2005} = \mathrm{i}^{2005} = (\mathrm{i}^4)^{501} \cdot \mathrm{i} = 1^{501} \cdot \mathrm{i} = \mathrm{i}.
\]
故选 A．

\gkanswer{答案：A}

\end{answer}
\begin{answer}
\gksolhead{3}{}

\textbf{题目}：

若函数 $f(x)$ 是定义在 $\mathbf{R}$ 上的偶函数，在 $(-\infty, 0]$ 上是减函数，且 $f(2) = 0$，则使得 $f(x) < 0$ 的 $x$ 的取值范围是\mbox{（　　）}

\begin{gkopts}(2)
\task A．$(-\infty, 2)$
\task B．$(2, +\infty)$
\task C．$(-\infty, -2) \cup (2, +\infty)$
\task D．$(-2, 2)$
\end{gkopts}

\Needspace{4\baselineskip}
\textbf{解析}

因为 $f(x)$ 是定义在 $\mathbf{R}$ 上的偶函数，其图象关于 $y$ 轴对称．
又 $f(x)$ 在 $(-\infty, 0]$ 上是减函数，由偶函数的对称性可知，$f(x)$ 在 $[0, +\infty)$ 上是增函数．
已知 $f(2) = 0$，由偶函数性质可得 $f(-2) = f(2) = 0$．
在区间 $[0, +\infty)$ 上：
当 $0 \leqslant x < 2$ 时，$f(x) < f(2) = 0$；
由偶函数图象的对称性，在区间 $(-\infty, 0]$ 上：
当 $-2 < x \leqslant 0$ 时，$f(x) < f(-2) = 0$．
综合两部分，使得 $f(x) < 0$ 的 $x$ 的取值范围是 $(-2, 2)$．
故选 D．

\gkanswer{答案：D}

\end{answer}
\begin{answer}
\gksolhead{4}{}

\textbf{题目}：

已知 $A(3, 1), B(6, 1), C(4, 3)$，$D$ 为线段 $BC$ 的中点，则向量 $\overrightarrow{AC}$ 与 $\overrightarrow{DA}$ 的夹角为\mbox{（　　）}

\begin{gkopts}(2)
\task A．$\dfrac{\pi}{2} - \arccos\dfrac{4}{5}$
\task B．$\arccos\dfrac{4}{5}$
\task C．$\arccos\left(-\dfrac{4}{5}\right)$
\task D．$-\arccos\left(-\dfrac{4}{5}\right)$
\end{gkopts}

\Needspace{4\baselineskip}
\textbf{解析}

\Needspace{5\baselineskip}
由已知 $B(6, 1), C(4, 3)$，且 $D$ 为线段 $BC$ 的中点，由中点坐标公式得：
\nopagebreak[4]
\[
D\left(\frac{6 + 4}{2}, \frac{1 + 3}{2}\right) = (5, 2).
\]
\Needspace{5\baselineskip}
又 $A(3, 1)$，计算相关向量的坐标：
\nopagebreak[4]
\[
\overrightarrow{AC} = (4 - 3, 3 - 1) = (1, 2),
\]
\nopagebreak[4]
\[
\overrightarrow{DA} = (3 - 5, 1 - 2) = (-2, -1).
\]
\Needspace{5\baselineskip}
计算两向量的模长和数量积：
\nopagebreak[4]
\[
|\overrightarrow{AC}| = \sqrt{1^2 + 2^2} = \sqrt{5},
\]
\nopagebreak[4]
\[
|\overrightarrow{DA}| = \sqrt{(-2)^2 + (-1)^2} = \sqrt{5},
\]
\nopagebreak[4]
\[
\overrightarrow{AC} \cdot \overrightarrow{DA} = 1 \times (-2) + 2 \times (-1) = -2 - 2 = -4.
\]
\Needspace{5\baselineskip}
设向量 $\overrightarrow{AC}$ 与 $\overrightarrow{DA}$ 的夹角为 $\theta$（$0 \leqslant \theta \leqslant \pi$），则：
\nopagebreak[4]
\[
\cos\theta = \frac{\overrightarrow{AC} \cdot \overrightarrow{DA}}{|\overrightarrow{AC}| |\overrightarrow{DA}|} = \frac{-4}{\sqrt{5} \times \sqrt{5}} = -\frac{4}{5}.
\]
因此 $\theta = \arccos\left(-\dfrac{4}{5}\right)$．
故选 C．

\gkanswer{答案：C}

\end{answer}
\begin{answer}
\gksolhead{5}{}

\textbf{题目}：

若 $x, y$ 是正数，则 $\left(x + \dfrac{1}{2y}\right)^2 + \left(y + \dfrac{1}{2x}\right)^2$ 的最小值是\mbox{（　　）}

\begin{gkopts}(4)
\task A．$3$
\task B．$\dfrac{7}{2}$
\task C．$4$
\task D．$\dfrac{9}{2}$
\end{gkopts}

\Needspace{4\baselineskip}
\textbf{解析}

\Needspace{5\baselineskip}
将原式展开变形：
\nopagebreak[4]
\[
\left(x + \frac{1}{2y}\right)^2 + \left(y + \frac{1}{2x}\right)^2 = x^2 + \frac{x}{y} + \frac{1}{4y^2} + y^2 + \frac{y}{x} + \frac{1}{4x^2}
\]
\nopagebreak[4]
\[
= \left(x^2 + \frac{1}{4x^2}\right) + \left(y^2 + \frac{1}{4y^2}\right) + \left(\frac{x}{y} + \frac{y}{x}\right).
\]
\Needspace{5\baselineskip}
因为 $x > 0, y > 0$，由基本不等式：
\nopagebreak[4]
\[
x^2 + \frac{1}{4x^2} \geqslant 2\sqrt{x^2 \cdot \frac{1}{4x^2}} = 1,
\]
当且仅当 $x^2 = \dfrac{1}{4x^2}$ 即 $x = \dfrac{\sqrt{2}}{2}$ 时等号成立；
\Needspace{5\baselineskip}
同理：
\nopagebreak[4]
\[
y^2 + \frac{1}{4y^2} \geqslant 2\sqrt{y^2 \cdot \frac{1}{4y^2}} = 1,
\]
当且仅当 $y = \dfrac{\sqrt{2}}{2}$ 时等号成立；
\Needspace{5\baselineskip}
又：
\nopagebreak[4]
\[
\frac{x}{y} + \frac{y}{x} \geqslant 2\sqrt{\frac{x}{y} \cdot \frac{y}{x}} = 2,
\]
当且仅当 $x = y$ 时等号成立．
当 $x = y = \dfrac{\sqrt{2}}{2}$ 时，上述三式等号同时成立，
因此原式的最小值为 $1 + 1 + 2 = 4$．
故选 C．

\gkanswer{答案：C}

\end{answer}
\begin{answer}
\gksolhead{6}{}

\textbf{题目}：

已知 $\alpha, \beta$ 均为锐角，若 $p: \sin\alpha < \sin(\alpha + \beta)$，$q: \alpha + \beta < \dfrac{\pi}{2}$，则 $p$ 是 $q$ 的\mbox{（　　）}

\begin{gkopts}(2)
\task A．充分而不必要条件
\task B．必要而不充分条件
\task C．充要条件
\task D．既不充分也不必要条件
\end{gkopts}

\Needspace{4\baselineskip}
\textbf{解析}

因为 $\alpha, \beta$ 均为锐角，所以 $0 < \alpha < \dfrac{\pi}{2}, 0 < \beta < \dfrac{\pi}{2}$，从而 $0 < \alpha + \beta < \pi$．
若条件 $q: \alpha + \beta < \dfrac{\pi}{2}$ 成立，
则 $\alpha$ 与 $\alpha + \beta$ 均属于区间 $\left(0, \dfrac{\pi}{2}\right)$．
因为正弦函数 $y = \sin x$ 在 $\left(0, \dfrac{\pi}{2}\right)$ 上单调递增，且 $\alpha < \alpha + \beta$，
所以 $\sin\alpha < \sin(\alpha + \beta)$，即命题 $p$ 成立．
这说明由 $q$ 可以推出 $p$．
反之，若条件 $p: \sin\alpha < \sin(\alpha + \beta)$ 成立，不能推出 $q: \alpha + \beta < \dfrac{\pi}{2}$．
例如取 $\alpha = \dfrac{\pi}{6}, \beta = \dfrac{\pi}{3}$（均为锐角），
此时 $\alpha + \beta = \dfrac{\pi}{2}$，$\sin\alpha = \dfrac{1}{2} < \sin(\alpha + \beta) = 1$，$p$ 成立，但 $q$ 不成立；
又如取 $\alpha = \dfrac{\pi}{6}, \beta = \dfrac{4\pi}{9}$，$\alpha + \beta = \dfrac{11\pi}{18} > \dfrac{\pi}{2}$，$\sin\dfrac{11\pi}{18} = \sin\dfrac{7\pi}{18} > \sin\dfrac{\pi}{6}$，$p$ 成立但 $q$ 不成立．
故由 $p$ 推不出 $q$．
综上可知，$p$ 是 $q$ 的必要而不充分条件．
故选 B．

\gkanswer{答案：B}

\end{answer}
\begin{answer}
\gksolhead{7}{}

\textbf{题目}：

对于不重合的两个平面 $\alpha$ 与 $\beta$，给定下列条件：

① 存在平面 $\gamma$，使得 $\alpha, \beta$ 都垂直于 $\gamma$；

② 存在平面 $\gamma$，使得 $\alpha, \beta$ 都平行于 $\gamma$；

③ $\alpha$ 内有不共线的三点到 $\beta$ 的距离相等；

④ 存在异面直线 $l, m$，使得 $l \parallel \alpha, l \parallel \beta, m \parallel \alpha, m \parallel \beta$．

其中，可以判定 $\alpha$ 与 $\beta$ 平行的条件有\mbox{（　　）}

\begin{gkopts}(4)
\task A．1 个
\task B．2 个
\task C．3 个
\task D．4 个
\end{gkopts}

\Needspace{4\baselineskip}
\textbf{解析}

逐项分析各条件：
① 存在平面 $\gamma$，使得 $\alpha \perp \gamma, \beta \perp \gamma$．垂直于同一个平面的两个不同平面可以平行也可以相交（例如教室内相邻两面墙都垂直于地面），故 ① 不能判定平行；
② 存在平面 $\gamma$，使得 $\alpha \parallel \gamma, \beta \parallel \gamma$．由面面平行的传递性，平行于同一平面的两个不同平面互相平行，故 ② 可以判定平行；
③ $\alpha$ 内有不共线的三点到 $\beta$ 的距离相等．若这三点分布在平面 $\beta$ 的同侧，则 $\alpha \parallel \beta$；若其中两点在 $\beta$ 的一侧，另一点在 $\beta$ 的另一侧且距离相等，则两平面相交，故 ③ 不能判定平行；
④ 存在异面直线 $l, m$，使得 $l \parallel \alpha, l \parallel \beta, m \parallel \alpha, m \parallel \beta$．
在平面 $\alpha$ 内任取一点 $P$，过 $P$ 分别作 $l' \parallel l, m' \parallel m$．
因为 $l$ 与 $m$ 是异面直线，所以 $l'$ 与 $m'$ 是平面 $\alpha$ 内的两条相交直线．
又因为 $l \parallel \beta, m \parallel \beta$，所以 $l' \parallel \beta, m' \parallel \beta$．
根据面面平行的判定定理：“一个平面内的两条相交直线与另一个平面平行，则这两个平面平行”，可得 $\alpha \parallel \beta$．故 ④ 可以判定平行．
综上，能够判定 $\alpha$ 与 $\beta$ 平行的条件有 ②④，共 2 个．
故选 B．

\gkanswer{答案：B}

\end{answer}
\begin{answer}
\gksolhead{8}{}

\textbf{题目}：

若 $\left(2x - \dfrac{1}{x}\right)^n$ 展开式中含 $\dfrac{1}{x^2}$ 项的系数与含 $\dfrac{1}{x^4}$ 项的系数之比为 $-5$，则 $n$ 等于\mbox{（　　）}

\begin{gkopts}(4)
\task A．4
\task B．6
\task C．8
\task D．10
\end{gkopts}

\Needspace{4\baselineskip}
\textbf{解析}

\Needspace{5\baselineskip}
二项式 $\left(2x - \dfrac{1}{x}\right)^n$ 展开式的通项公式为：
\nopagebreak[4]
\[
T_{r+1} = \mathrm{C}_n^r (2x)^{n-r} \left(-\frac{1}{x}\right)^r = (-1)^r \mathrm{C}_n^r 2^{n-r} x^{n - 2r}.
\]
\Needspace{5\baselineskip}
令 $n - 2r = -2$，解得 $r = \dfrac{n+2}{2}$，此时对应含 $\dfrac{1}{x^2}$ 项的系数为：
\nopagebreak[4]
\[
A = (-1)^{\frac{n+2}{2}} \mathrm{C}_n^{\frac{n+2}{2}} 2^{n - \frac{n+2}{2}} = (-1)^{\frac{n+2}{2}} \mathrm{C}_n^{\frac{n+2}{2}} 2^{\frac{n-2}{2}}.
\]
\Needspace{5\baselineskip}
令 $n - 2r = -4$，解得 $r = \dfrac{n+4}{2}$，此时对应含 $\dfrac{1}{x^4}$ 项的系数为：
\nopagebreak[4]
\[
B = (-1)^{\frac{n+4}{2}} \mathrm{C}_n^{\frac{n+4}{2}} 2^{n - \frac{n+4}{2}} = (-1)^{\frac{n+4}{2}} \mathrm{C}_n^{\frac{n+4}{2}} 2^{\frac{n-4}{2}}.
\]
\Needspace{5\baselineskip}
两系数之比为：
\nopagebreak[4]
\[
\frac{A}{B} = \frac{(-1)^{\frac{n+2}{2}} \mathrm{C}_n^{\frac{n+2}{2}} 2^{\frac{n-2}{2}}}{(-1)^{\frac{n+4}{2}} \mathrm{C}_n^{\frac{n+4}{2}} 2^{\frac{n-4}{2}}} = -\frac{2 \mathrm{C}_n^{\frac{n+2}{2}}}{\mathrm{C}_n^{\frac{n+4}{2}}}.
\]
\Needspace{5\baselineskip}
已知 $\dfrac{A}{B} = -5$，故：
\nopagebreak[4]
\[
\frac{2 \mathrm{C}_n^{\frac{n+2}{2}}}{\mathrm{C}_n^{\frac{n+4}{2}}} = 5.
\]
\Needspace{5\baselineskip}
记 $k = \dfrac{n+2}{2}$，则 $\dfrac{n+4}{2} = k + 1$．由组合数性质：
\nopagebreak[4]
\[
\frac{\mathrm{C}_n^k}{\mathrm{C}_n^{k+1}} = \frac{k+1}{n-k}.
\]
\Needspace{5\baselineskip}
代入 $k = \dfrac{n+2}{2}$，得：
\nopagebreak[4]
\[
\frac{k+1}{n-k} = \frac{\frac{n+4}{2}}{\frac{n-2}{2}} = \frac{n+4}{n-2}.
\]
\Needspace{5\baselineskip}
因此：
\nopagebreak[4]
\[
2 \cdot \frac{n+4}{n-2} = 5 \implies 2n + 8 = 5n - 10 \implies 3n = 18 \implies n = 6.
\]
故选 B．

\gkanswer{答案：B}

\end{answer}
\begin{answer}
\gksolhead{9}{}

\textbf{题目}：

若动点 $(x, y)$ 在曲线 $\dfrac{x^2}{4} + \dfrac{y^2}{b^2} = 1$ ($b > 0$) 上变化，则 $x^2 + 2y$ 的最大值为\mbox{（　　）}

\begin{gkopts}(1)
\task A．$\begin{cases} \dfrac{b^2}{4} + 4, & 0 < b < 4, \\ 2b, & b \geqslant 4. \end{cases}$
\task B．$\begin{cases} \dfrac{b^2}{4} + 4, & 0 < b < 2, \\ 2b, & b \geqslant 2. \end{cases}$
\task C．$\dfrac{b^2}{4} + 4$
\task D．$2b$
\end{gkopts}

\Needspace{4\baselineskip}
\textbf{解析}

由动点 $(x, y)$ 在椭圆 $\dfrac{x^2}{4} + \dfrac{y^2}{b^2} = 1$ ($b > 0$) 上变化，
可得 $x^2 = 4\left(1 - \dfrac{y^2}{b^2}\right) = 4 - \dfrac{4}{b^2}y^2$，其中 $y \in [-b, b]$．
\Needspace{5\baselineskip}
将其代入目标代数式 $x^2 + 2y$，设函数：
\nopagebreak[4]
\[
g(y) = 4 - \frac{4}{b^2}y^2 + 2y = -\frac{4}{b^2}\left(y^2 - \frac{b^2}{2}y\right) + 4
\]
\nopagebreak[4]
\[
= -\frac{4}{b^2}\left(y - \frac{b^2}{4}\right)^2 + \frac{b^2}{4} + 4.
\]
该二次函数图象开口向下，对称轴方程为 $y = \dfrac{b^2}{4}$．
由于定义域为 $y \in [-b, b]$，需根据对称轴与区间右端点 $b$ 的相对位置分类讨论：
1. 当 $\dfrac{b^2}{4} < b$ 即 $0 < b < 4$ 时：
对称轴 $y = \dfrac{b^2}{4} \in (-b, b)$ 在定义域内部，
因此当 $y = \dfrac{b^2}{4}$ 时，函数取得最大值 $\dfrac{b^2}{4} + 4$；
2. 当 $\dfrac{b^2}{4} \geqslant b$ 即 $b \geqslant 4$ 时：
对称轴在区间 $[-b, b]$ 右侧（或与右端点重合），
此时函数 $g(y)$ 在 $[-b, b]$ 上单调递增，最大值在右端点 $y = b$ 处取得，
最大值为 $g(b) = -\dfrac{4}{b^2}(b^2) + 2b + 4 = -4 + 2b + 4 = 2b$．
\Needspace{5\baselineskip}
综合上述讨论，最大值为：
\nopagebreak[4]
\[
\begin{cases}
\dfrac{b^2}{4} + 4, & 0 < b < 4, \\
2b, & b \geqslant 4.
\end{cases}
\]
故选 A．

\gkanswer{答案：A}

\end{answer}
\begin{answer}
\gksolhead{10}{}

\textbf{题目}：

在体积为 1 的三棱锥 $A - BCD$ 侧棱 $AB, AC, AD$ 上分别取点 $E, F, G$，使 $AE:EB = AF:FC = AG:GD = 2:1$，记 $O$ 为三平面 $BCG, CDE, DBF$ 的交点，则三棱锥 $O - BCD$ 的体积等于\mbox{（　　）}

\begin{gkopts}(4)
\task A．$\dfrac{1}{9}$
\task B．$\dfrac{1}{8}$
\task C．$\dfrac{1}{7}$
\task D．$\dfrac{1}{4}$
\end{gkopts}

\Needspace{4\baselineskip}
\textbf{解析}

设三棱锥 $A - BCD$ 的体积为 $V = 1$．
\Needspace{5\baselineskip}
因为 $E, F, G$ 分别在侧棱 $AB, AC, AD$ 上，且满足：
\nopagebreak[4]
\[
AE : EB = AF : FC = AG : GD = 2 : 1,
\]
所以 $\dfrac{AE}{AB} = \dfrac{AF}{AC} = \dfrac{AG}{AD} = \dfrac{2}{3}$，平面 $EFG \parallel$ 底面 $BCD$．
引入四面体质点系（或重心坐标）：
设顶点 $A, B, C, D$ 分别分配质点权重 $w_A, w_B, w_C, w_D$．
点 $E$ 在线段 $AB$ 上，满足 $AE : EB = 2 : 1$．由杠杆原理，点 $E$ 是顶点 $A$ 处权重为 1、顶点 $B$ 处权重为 2 的质点系的质心．
平面 $CDE$ 经过直线 $CD$ 和点 $E$，故空间点 $X$ 在平面 $CDE$ 上当且仅当其在 $\{A, B, C, D\}$ 下的重心坐标满足 $x_B = 2x_A$．
同理：
点 $F$ 在线段 $AC$ 上，满足 $AF : FC = 2 : 1$，故点 $X$ 在平面 $DBF$ 上当且仅当 $x_C = 2x_A$；
点 $G$ 在线段 $AD$ 上，满足 $AG : GD = 2 : 1$，故点 $X$ 在平面 $BCG$ 上当且仅当 $x_D = 2x_A$．
\Needspace{5\baselineskip}
点 $O$ 为三平面 $BCG, CDE, DBF$ 的公共交点，因此点 $O$ 的重心坐标必须同时满足：
\nopagebreak[4]
\[
x_B = 2x_A, \quad x_C = 2x_A, \quad x_D = 2x_A.
\]
\Needspace{5\baselineskip}
由重心坐标的归一化条件：
\nopagebreak[4]
\[
x_A + x_B + x_C + x_D = 1 \implies x_A + 2x_A + 2x_A + 2x_A = 1 \implies 7x_A = 1 \implies x_A = \frac{1}{7}.
\]
因为点 $O$ 到底面 $BCD$ 的距离占顶点 $A$ 到底面 $BCD$ 距离的比例恰好等于重心坐标 $x_A$，
所以三棱锥 $O - BCD$ 的高是三棱锥 $A - BCD$ 的高的 $\dfrac{1}{7}$．
\Needspace{5\baselineskip}
因为两三棱锥同底（底面均为 $\triangle BCD$），所以：
\nopagebreak[4]
\[
V_{O-BCD} = x_A \cdot V_{A-BCD} = \frac{1}{7} \times 1 = \frac{1}{7}.
\]
故选 C．

\gkanswer{答案：C}

\end{answer}
\begin{answer}
\gksolhead{11}{}

\textbf{题目}：

已知集合 $A = \{x \in \mathbf{R} \mid x^2 - x - 6 < 0\}$，集合 $B = \{x \in \mathbf{R} \mid |x - 2| < 2\}$，则 $A \cap B =$ \gkblank{}．

\Needspace{4\baselineskip}
\textbf{解析}

\Needspace{5\baselineskip}
解不等式 $x^2 - x - 6 < 0$：
\nopagebreak[4]
\[
(x - 3)(x + 2) < 0 \iff -2 < x < 3,
\]
所以集合 $A = (-2, 3)$．
\Needspace{5\baselineskip}
解不等式 $|x - 2| < 2$：
\nopagebreak[4]
\[
-2 < x - 2 < 2 \iff 0 < x < 4,
\]
所以集合 $B = (0, 4)$．
\Needspace{5\baselineskip}
求两集合的交集：
\nopagebreak[4]
\[
A \cap B = (-2, 3) \cap (0, 4) = (0, 3).
\]
故填 $(0, 3)$（或 $\{x \mid 0 < x < 3\}$）．

\gkanswer{答案：$(0, 3)$（或 $\{x \mid 0 < x < 3\}$）}

\end{answer}
\begin{answer}
\gksolhead{12}{}

\textbf{题目}：

曲线 $y = x^3$ 在点 $(a, a^3)$ ($a \neq 0$) 处的切线与 $x$ 轴、直线 $x = a$ 所围成的三角形的面积为 $\dfrac{1}{6}$，则 $a =$ \gkblank{}．

\Needspace{4\baselineskip}
\textbf{解析}

函数 $y = x^3$ 的导数为 $y' = 3x^2$．
在点 $(a, a^3)$ ($a \neq 0$) 处，切线斜率 $k = 3a^2$．
\Needspace{5\baselineskip}
切线方程为：
\nopagebreak[4]
\[
y - a^3 = 3a^2(x - a) \implies y = 3a^2 x - 2a^3.
\]
求切线与 $x$ 轴的交点：令 $y = 0$，得 $3a^2 x = 2a^3$．
因为 $a \neq 0$，解得 $x = \dfrac{2}{3}a$，即交点坐标为 $\left(\dfrac{2}{3}a, 0\right)$．
直线 $x = a$ 与 $x$ 轴的交点坐标为 $(a, 0)$，与切线的交点为切点 $(a, a^3)$．
切线与 $x$ 轴、直线 $x = a$ 所围成的图形是以 $\left(\dfrac{2}{3}a, 0\right)$、$(a, 0)$、$(a, a^3)$ 为顶点的直角三角形，
其在 $x$ 轴上的底边长为 $\left|a - \dfrac{2}{3}a\right| = \dfrac{1}{3}|a|$，
高为线段长 $|a^3|$．
\Needspace{5\baselineskip}
因此三角形的面积为：
\nopagebreak[4]
\[
S = \frac{1}{2} \times \frac{1}{3}|a| \times |a^3| = \frac{1}{6} a^4.
\]
\Needspace{5\baselineskip}
已知面积为 $\dfrac{1}{6}$，所以：
\nopagebreak[4]
\[
\frac{1}{6} a^4 = \frac{1}{6} \implies a^4 = 1 \implies a = \pm 1.
\]
故填 $\pm 1$（或 $1$ 或 $-1$）．

\gkanswer{答案：$\pm 1$（或 $1$ 或 $-1$）}

\end{answer}
\begin{answer}
\gksolhead{13}{}

\textbf{题目}：

已知 $\alpha, \beta$ 均为锐角，且 $\cos(\alpha + \beta) = \sin(\alpha - \beta)$，则 $\tan\alpha =$ \gkblank{}．

\Needspace{4\baselineskip}
\textbf{解析}

\Needspace{5\baselineskip}
利用三角函数的两角和与差展开公式，由已知：
\nopagebreak[4]
\[
\cos(\alpha + \beta) = \sin(\alpha - \beta),
\]
\Needspace{5\baselineskip}
展开得：
\nopagebreak[4]
\[
\cos\alpha\cos\beta - \sin\alpha\sin\beta = \sin\alpha\cos\beta - \cos\alpha\sin\beta.
\]
\Needspace{5\baselineskip}
移项整理：
\nopagebreak[4]
\[
\cos\alpha\cos\beta + \cos\alpha\sin\beta = \sin\alpha\cos\beta + \sin\alpha\sin\beta
\]
\nopagebreak[4]
\[
\implies \cos\alpha(\cos\beta + \sin\beta) = \sin\alpha(\cos\beta + \sin\beta).
\]
因为 $\beta$ 为锐角，所以 $\cos\beta > 0, \sin\beta > 0$，从而 $\cos\beta + \sin\beta > 0$．
\Needspace{5\baselineskip}
在等式两边同除以正数 $\cos\beta + \sin\beta$，得：
\nopagebreak[4]
\[
\cos\alpha = \sin\alpha.
\]
\Needspace{5\baselineskip}
又因为 $\alpha$ 为锐角，$\cos\alpha > 0$，两边同除以 $\cos\alpha$，得：
\nopagebreak[4]
\[
\tan\alpha = 1.
\]
故填 $1$．

\gkanswer{答案：$1$}

\end{answer}
\begin{answer}
\gksolhead{14}{}

\textbf{题目}：

$\lim_{n \to \infty} \dfrac{2^{3n} - 3^{2n+1}}{2^{3n} + 3^{2n}} =$ \gkblank{}．

\Needspace{4\baselineskip}
\textbf{解析}

\Needspace{5\baselineskip}
将分子与分母中的指数式化为以同底数为基准：
\nopagebreak[4]
\[
2^{3n} = (2^3)^n = 8^n, \quad 3^{2n} = (3^2)^n = 9^n, \quad 3^{2n+1} = 3 \cdot 9^n.
\]
\Needspace{5\baselineskip}
代入原式：
\nopagebreak[4]
\[
\lim_{n \to \infty} \frac{2^{3n} - 3^{2n+1}}{2^{3n} + 3^{2n}} = \lim_{n \to \infty} \frac{8^n - 3 \cdot 9^n}{8^n + 9^n}.
\]
\Needspace{5\baselineskip}
分子分母同除以最大底数项 $9^n$：
\nopagebreak[4]
\[
= \lim_{n \to \infty} \frac{\left(\dfrac{8}{9}\right)^n - 3}{\left(\dfrac{8}{9}\right)^n + 1}.
\]
因为 $0 < \dfrac{8}{9} < 1$，当 $n \to \infty$ 时，$\lim_{n \to \infty} \left(\dfrac{8}{9}\right)^n = 0$．
\Needspace{5\baselineskip}
所以原极限为：
\nopagebreak[4]
\[
\frac{0 - 3}{0 + 1} = -3.
\]
故填 $-3$．

\gkanswer{答案：$-3$}

\end{answer}
\begin{answer}
\gksolhead{15}{}

\textbf{题目}：

某轻轨列车有 4 节车厢，现有 6 位乘客准备乘坐，设每一位乘客进入每节车厢是等可能的，则这 6 位乘客进入各节车厢的人数恰好为 0, 1, 2, 3 的概率为 \gkblank{}．

\Needspace{4\baselineskip}
\textbf{解析}

\Needspace{5\baselineskip}
每位乘客进入每节车厢有 4 种等可能选择，6 位乘客进入 4 节车厢的全部可能情况总数为：
\nopagebreak[4]
\[
N = 4^6 = 4096.
\]
求 6 位乘客进入各节车厢的人数恰好为 0, 1, 2, 3 的情况数：
第一步，将 4 节车厢分别指定容纳 0, 1, 2, 3 位乘客，共有 $\mathrm{A}_4^4 = 24$ 种分配方案；
第二步，将 6 位乘客分成四组，人数分别为 0, 1, 2, 3：
\begin{itemize}
\item 容纳 0 人的车厢无乘客，方案数为 $\mathrm{C}_6^0 = 1$；
\item 从 6 位乘客中选取 1 人进入容纳 1 人的车厢，方案数为 $\mathrm{C}_6^1$；
\item 从剩余 5 位乘客中选取 2 人进入容纳 2 人的车厢，方案数为 $\mathrm{C}_5^2$；
\item 剩余 3 位乘客全部进入容纳 3 人的车厢，方案数为 $\mathrm{C}_3^3 = 1$．
\end{itemize}
分组方案数为 $\mathrm{C}_6^1 \mathrm{C}_5^2 = 6 \times 10 = 60$．
\Needspace{5\baselineskip}
根据分步乘法计数原理，满足条件的基本事件总数为：
\nopagebreak[4]
\[
M = \mathrm{A}_4^4 \times \mathrm{C}_6^1 \mathrm{C}_5^2 = 24 \times 60 = 1440.
\]
\Needspace{5\baselineskip}
因此所求概率为：
\nopagebreak[4]
\[
P = \frac{M}{N} = \frac{1440}{4096} = \frac{45}{128}.
\]
故填 $\dfrac{45}{128}$．

\gkanswer{答案：$\dfrac{45}{128}$}

\end{answer}
\begin{answer}
\gksolhead{16}{}

\textbf{题目}：

连接抛物线上任意四点组成的四边形可能是 \gkblank{}．（填写所有正确选项的序号）

① 菱形；② 有 3 条边相等的四边形；③ 梯形；④ 平行四边形；⑤ 有一组对角相等的四边形．

\Needspace{4\baselineskip}
\textbf{解析}

设抛物线标准方程为 $y^2 = 2px$ ($p > 0$)，其为严格凸曲线．
1. \textbf{关于平行四边形与菱形（①、④）}：
假设抛物线上存在四点 $A, B, C, D$ 构成平行四边形，则必有 $\overrightarrow{AB} = \overrightarrow{DC}$．
设四点坐标为 $(x_k, y_k)$ ($k = 1, 2, 3, 4$)，且 $x_k = \dfrac{y_k^2}{2p}$．
由 $\overrightarrow{AB} = \overrightarrow{DC}$，纵坐标满足 $y_2 - y_1 = y_3 - y_4$；
\Needspace{5\baselineskip}
横坐标满足 $x_2 - x_1 = x_3 - x_4$，即：
\nopagebreak[4]
\[
\frac{y_2^2 - y_1^2}{2p} = \frac{y_3^2 - y_4^2}{2p} \implies (y_2 - y_1)(y_2 + y_1) = (y_3 - y_4)(y_3 + y_4).
\]
若 $y_2 - y_1 = y_3 - y_4 \neq 0$，则必有 $y_2 + y_1 = y_3 + y_4$．
联立两式相加可得 $2y_2 = 2y_3 \implies y_2 = y_3$，从而 $B$ 与 $C$ 重合，与四边形顶点互异矛盾；
若 $y_2 - y_1 = 0$，则直线 $AB$ 平行于对称轴（$x$ 轴），但平行于对称轴的直线与抛物线至多只有 1 个交点，不可能包含 2 个不同的顶点．
因此抛物线上任意四点不可能组成平行四边形，更不可能组成菱形．故 ①、④ 错误；
2. \textbf{关于梯形（③）}：
取关于抛物线对称轴对称的四点，例如横坐标为 $x_1$ 处的两点 $(x_1, \pm y_1)$ 与横坐标为 $x_2$ 处的两点 $(x_2, \pm y_2)$（$x_1 \neq x_2$）．
这两对点连接构成的两条底边垂直于对称轴且平行，两底长度不等，连接对应端点即构成等腰梯形．故 ③ 正确；
3. \textbf{关于有 3 条边相等的四边形（②）}：
在上述等腰梯形构型中，两腰长始终相等，通过连续变动端点坐标，可以使得上底长等于两腰长，从而构成有 3 条边相等的四边形．故 ② 正确；
4. \textbf{关于有一组对角相等的四边形（⑤）}：
抛物线上四点确定一个四边形共有 4 个自由度，而“一组对角相等”是一个一维等式约束．在一般非对称或一般凸四边形微调中，该约束显然可以被满足．故 ⑤ 正确．
综上所述，所有正确选项的序号为 ②③⑤．
故填 ②③⑤．

\gkanswer{答案：②③⑤}

\end{answer}
\begin{answer}
\gksolhead{17}{}

\textbf{题目}：

（本小题满分 12 分）

若函数 $f(x) = \dfrac{1 + \cos 2x}{4\sin\left(\dfrac{\pi}{2} + x\right)} - a \sin\dfrac{x}{2} \cos\left(\pi - \dfrac{x}{2}\right)$ 的最大值为 2，试确定常数 $a$ 的值．

\Needspace{4\baselineskip}
\textbf{解析}

首先对函数 $f(x)$ 的解析式进行三角恒等化简：
因为 $1 + \cos 2x = 2\cos^2 x$，
且 $\sin\left(\dfrac{\pi}{2} + x\right) = \cos x$，
\Needspace{5\baselineskip}
所以当 $\cos x \neq 0$ 时：
\nopagebreak[4]
\[
\frac{1 + \cos 2x}{4\sin\left(\dfrac{\pi}{2} + x\right)} = \frac{2\cos^2 x}{4\cos x} = \frac{1}{2}\cos x.
\]
又因为 $\cos\left(\pi - \dfrac{x}{2}\right) = -\cos\dfrac{x}{2}$，
\Needspace{5\baselineskip}
所以第二项为：
\nopagebreak[4]
\[
-a \sin\frac{x}{2} \cos\left(\pi - \frac{x}{2}\right) = a \sin\frac{x}{2}\cos\frac{x}{2} = \frac{a}{2}\sin x.
\]
\Needspace{5\baselineskip}
因此：
\nopagebreak[4]
\[
f(x) = \frac{1}{2}\cos x + \frac{a}{2}\sin x = \frac{1}{2}(\cos x + a \sin x).
\]
\Needspace{5\baselineskip}
引入辅助角 $\varphi$（其中 $\tan\varphi = \dfrac{1}{a}$），可化为：
\nopagebreak[4]
\[
f(x) = \frac{\sqrt{a^2 + 1}}{2} \sin(x + \varphi).
\]
因为当 $\sin(x + \varphi) = 1$ 时，$\cos x = \sin\varphi = \dfrac{1}{\sqrt{a^2 + 1}} \neq 0$，最值点满足分母不为 0 的定义域要求，
所以函数 $f(x)$ 的最大值为 $\dfrac{\sqrt{a^2 + 1}}{2}$．
\Needspace{5\baselineskip}
由题意，$f(x)$ 的最大值为 2，故：
\nopagebreak[4]
\[
\frac{\sqrt{a^2 + 1}}{2} = 2 \implies \sqrt{a^2 + 1} = 4 \implies a^2 + 1 = 16 \implies a^2 = 15.
\]
解得 $a = \pm\sqrt{15}$．

\gkanswer{答案：$a = \pm\sqrt{15}$。}

\end{answer}
\begin{answer}
\gksolhead{18}{}

\textbf{题目}：

（本小题满分 12 分）

在一次购物抽奖活动中，假设某 10 张券中有一等奖券 1 张，可获价值 50 元的奖品；有二等奖券 3 张，每张可获价值 10 元的奖品；其余 6 张没有奖，某顾客从此 10 张券中任抽 2 张，求：

(1) 该顾客中奖的概率；

(2) 该顾客获得的奖品总价值 $\xi$（元）的概率分布列和期望 $\mathrm{E}\xi$．

\Needspace{4\baselineskip}
\textbf{解析}

\Needspace{5\baselineskip}
设从 10 张券中任抽 2 张的全部基本事件总数为：
\nopagebreak[4]
\[
N = \mathrm{C}_{10}^2 = \frac{10 \times 9}{2} = 45.
\]

\gksubpart{(1)}

设事件 $A$ 为“该顾客中奖”，则对立事件 $\overline{A}$ 为“抽取的 2 张券均没有奖”．
\Needspace{5\baselineskip}
因为 10 张券中没有奖的券共有 6 张，所以抽取的 2 张均无奖的方法数为：
\nopagebreak[4]
\[
\mathrm{C}_6^2 = \frac{6 \times 5}{2} = 15.
\]
\Needspace{5\baselineskip}
故抽取的 2 张均无奖的概率为：
\nopagebreak[4]
\[
P(\overline{A}) = \frac{\mathrm{C}_6^2}{\mathrm{C}_{10}^2} = \frac{15}{45} = \frac{1}{3}.
\]
\Needspace{5\baselineskip}
因此该顾客中奖的概率为：
\nopagebreak[4]
\[
P(A) = 1 - P(\overline{A}) = 1 - \frac{1}{3} = \frac{2}{3}.
\]

\gksubpart{(2)}

顾客获得的奖品总价值 $\xi$ 的可能取值为 $0, 10, 20, 50, 60$（元）．
分别计算各取值的概率：
\begin{itemize}
\item 当 $\xi = 0$ 时（抽得 2 张无奖券）：
\end{itemize}
\nopagebreak[4]
\[
P(\xi = 0) = \frac{\mathrm{C}_6^2}{\mathrm{C}_{10}^2} = \frac{15}{45} = \frac{1}{3};
\]
\begin{itemize}
\item 当 $\xi = 10$ 时（抽得 1 张二等奖券和 1 张无奖券）：
\end{itemize}
\nopagebreak[4]
\[
P(\xi = 10) = \frac{\mathrm{C}_3^1 \mathrm{C}_6^1}{\mathrm{C}_{10}^2} = \frac{3 \times 6}{45} = \frac{18}{45} = \frac{2}{5};
\]
\begin{itemize}
\item 当 $\xi = 20$ 时（抽得 2 张二等奖券）：
\end{itemize}
\nopagebreak[4]
\[
P(\xi = 20) = \frac{\mathrm{C}_3^2}{\mathrm{C}_{10}^2} = \frac{3}{45} = \frac{1}{15};
\]
\begin{itemize}
\item 当 $\xi = 50$ 时（抽得 1 张一等奖券和 1 张无奖券）：
\end{itemize}
\nopagebreak[4]
\[
P(\xi = 50) = \frac{\mathrm{C}_1^1 \mathrm{C}_6^1}{\mathrm{C}_{10}^2} = \frac{1 \times 6}{45} = \frac{6}{45} = \frac{2}{15};
\]
\begin{itemize}
\item 当 $\xi = 60$ 时（抽得 1 张一等奖券和 1 张二等奖券）：
\end{itemize}
\nopagebreak[4]
\[
P(\xi = 60) = \frac{\mathrm{C}_1^1 \mathrm{C}_3^1}{\mathrm{C}_{10}^2} = \frac{1 \times 3}{45} = \frac{3}{45} = \frac{1}{15}.
\]
故 $\xi$ 的概率分布列为：

\begin{center}
\begin{tabular}{cccccc}
\toprule
$\xi$ & 0 & 10 & 20 & 50 & 60 \\
\midrule
$P$ & $\dfrac{1}{3}$ & $\dfrac{2}{5}$ & $\dfrac{1}{15}$ & $\dfrac{2}{15}$ & $\dfrac{1}{15}$ \\
\bottomrule
\end{tabular}
\end{center}

\Needspace{5\baselineskip}
计算 $\xi$ 的数学期望：
\nopagebreak[4]
\[
\mathrm{E}\xi = 0 \times \frac{15}{45} + 10 \times \frac{18}{45} + 20 \times \frac{3}{45} + 50 \times \frac{6}{45} + 60 \times \frac{3}{45}
\]
\nopagebreak[4]
\[
= \frac{0 + 180 + 60 + 300 + 180}{45} = \frac{720}{45} = 16\text{（\text{元}）}.
\]

\gkanswer{答案（1）：$\dfrac{2}{3}$。}

\gkanswer{答案（2）：$\xi$ 的分布列为 $P(\xi=0)=\dfrac{1}{3}, P(\xi=10)=\dfrac{2}{5}, P(\xi=20)=\dfrac{1}{15}, P(\xi=50)=\dfrac{2}{15}, P(\xi=60)=\dfrac{1}{15}$；数学期望 $\mathrm{E}\xi = 16$（元）。}

\end{answer}
\begin{answer}
\gksolhead{19}{}

\textbf{题目}：

（本小题满分 12 分）

已知 $a \in \mathbf{R}$，讨论函数 $f(x) = \mathrm{e}^x(x^2 + ax + a + 1)$ 的极值点的个数．

\Needspace{4\baselineskip}
\textbf{解析}

函数 $f(x) = \mathrm{e}^x(x^2 + ax + a + 1)$ 的定义域为 $\mathbf{R}$．
\Needspace{5\baselineskip}
对 $f(x)$ 求导：
\nopagebreak[4]
\[
f'(x) = (\mathrm{e}^x)'(x^2 + ax + a + 1) + \mathrm{e}^x(x^2 + ax + a + 1)'
\]
\nopagebreak[4]
\[
= \mathrm{e}^x (x^2 + ax + a + 1) + \mathrm{e}^x (2x + a) = \mathrm{e}^x [x^2 + (a + 2)x + 2a + 1].
\]
因为对任意 $x \in \mathbf{R}$ 恒有 $\mathrm{e}^x > 0$，所以 $f'(x)$ 的符号由二次多项式 $g(x) = x^2 + (a + 2)x + 2a + 1$ 决定．
\Needspace{5\baselineskip}
计算方程 $g(x) = 0$ 的判别式：
\nopagebreak[4]
\[
\Delta = (a + 2)^2 - 4(2a + 1) = a^2 + 4a + 4 - 8a - 4 = a^2 - 4a = a(a - 4).
\]
根据判别式 $\Delta$ 的符号分类讨论如下：
1. 当 $\Delta < 0$，即 $0 < a < 4$ 时：
$g(x) > 0$ 在 $\mathbf{R}$ 上恒成立，从而 $f'(x) > 0$ 在 $\mathbf{R}$ 上恒成立．
此时函数 $f(x)$ 在 $\mathbf{R}$ 上单调递增，没有极值点，极值点个数为 0；
2. 当 $\Delta = 0$，即 $a = 0$ 或 $a = 4$ 时：
\begin{itemize}
\item 若 $a = 0$，$g(x) = x^2 + 2x + 1 = (x + 1)^2 \geqslant 0$；
\item 若 $a = 4$，$g(x) = x^2 + 6x + 9 = (x + 3)^2 \geqslant 0$．
\end{itemize}
在这两种情况下，$f'(x) \geqslant 0$ 在 $\mathbf{R}$ 上恒成立，且仅在孤立点处 $f'(x) = 0$，导数在零点两侧不改变符号，函数 $f(x)$ 在 $\mathbf{R}$ 上仍为单调递增函数，没有极值点，极值点个数为 0；
3. 当 $\Delta > 0$，即 $a < 0$ 或 $a > 4$ 时：
方程 $g(x) = 0$ 有两个不相等的实数根，设为 $x_1, x_2$（不妨设 $x_1 < x_2$）．
当 $x \in (-\infty, x_1)$ 时，$g(x) > 0 \implies f'(x) > 0$，$f(x)$ 单调递增；
当 $x \in (x_1, x_2)$ 时，$g(x) < 0 \implies f'(x) < 0$，$f(x)$ 单调递减；
当 $x \in (x_2, +\infty)$ 时，$g(x) > 0 \implies f'(x) > 0$，$f(x)$ 单调递增．
因此 $x = x_1$ 为函数 $f(x)$ 的极大值点，$x = x_2$ 为极小值点，极值点个数为 2．
综上所述：
\begin{itemize}
\item 当 $0 \leqslant a \leqslant 4$ 时，$f(x)$ 的极值点个数为 0；
\item 当 $a < 0$ 或 $a > 4$ 时，$f(x)$ 的极值点个数为 2．
\end{itemize}

\gkanswer{答案：当 $0 \leqslant a \leqslant 4$ 时，函数 $f(x)$ 的极值点个数为 0；当 $a < 0$ 或 $a > 4$ 时，函数 $f(x)$ 的极值点个数为 2。}

\end{answer}
\begin{answer}
\gksolhead{20}{}

\textbf{题目}：

（本小题满分 12 分）

如图，在三棱柱 $ABC - A_1B_1C_1$ 中，$AB \perp$ 侧面 $BB_1C_1C$，$E$ 为棱 $CC_1$ 上异于 $C, C_1$ 的一点，$EA \perp EB_1$．已知 $AB = \sqrt{2}, BB_1 = 2, BC = 1, \angle BCC_1 = \dfrac{\pi}{3}$，求：

(1) 异面直线 $AB$ 与 $EB_1$ 的距离；

(2) 二面角 $A - EB_1 - A_1$ 的平面角的正切值．
\nopagebreak[4]
\begin{center}\begin{tikzpicture}[>=stealth, scale=1.3, line join=round, line cap=round, thick]
  % Coordinates
  \coordinate (A) at (0, 2.5);
  \coordinate (A1) at (3.5, 2.5);
  \coordinate (B) at (0, 0.6);
  \coordinate (B1) at (3.5, 0.6);
  \coordinate (C) at (-1.1, 0.2);
  \coordinate (C1) at (2.4, 0.2);
  \coordinate (E) at ($(C)!0.5!(C1)$);

  % Dashed lines (hidden / interior)
  \draw[dashed] (A) -- (B);
  \draw[dashed] (B) -- (C);
  \draw[dashed] (B) -- (B1);
  \draw[dashed] (B) -- (E);
  \draw[dashed] (E) -- (B1);
  \draw[dashed] (A) -- (B1);

  % Solid lines (outer / visible)
  \draw (A) -- (A1) -- (B1);
  \draw (A) -- (C) -- (C1) -- (B1);
  \draw (A1) -- (C1);
  \draw (A) -- (E);
  \draw (A1) -- (E);

  % Labels
  \node[above left] at (A) {$A$};
  \node[above right] at (A1) {$A_1$};
  \node[above left] at (B) {$B$};
  \node[right] at (B1) {$B_1$};
  \node[below left] at (C) {$C$};
  \node[below] at (E) {$E$};
  \node[below right] at (C1) {$C_1$};
\end{tikzpicture}
\end{center}
\nopagebreak[4]
\Needspace{4\baselineskip}
\textbf{解析}

\gksubpart{(1)}

求解异面直线 $AB$ 与 $EB_1$ 的距离：
由已知 $AB \perp$ 侧面 $BB_1C_1C$，$B$ 为垂足．
直线 $EB_1$ 位于侧面 $BB_1C_1C$ 内，故 $AB \perp EB_1$．
在侧面 $BB_1C_1C$ 内，侧面为平行四边形，$BC = 1, BB_1 = CC_1 = 2, \angle BCC_1 = 60^\circ$．
因为 $E$ 在棱 $CC_1$ 上，设 $CE = x$（$0 < x < 2$），则 $EC_1 = 2 - x$．
\Needspace{5\baselineskip}
在 $\triangle BCE$ 中，由余弦定理：
\nopagebreak[4]
\[
BE^2 = BC^2 + CE^2 - 2 BC \cdot CE \cos 60^\circ = 1 + x^2 - x.
\]
\Needspace{5\baselineskip}
在 $\triangle B_1C_1E$ 中，因为 $\angle B_1C_1C = 180^\circ - \angle BCC_1 = 120^\circ$，由余弦定理：
\nopagebreak[4]
\[
EB_1^2 = B_1C_1^2 + EC_1^2 - 2 B_1C_1 \cdot EC_1 \cos 120^\circ = 1 + (2 - x)^2 - 2(1)(2 - x)\left(-\frac{1}{2}\right) = x^2 - 5x + 7.
\]
又因为 $AB \perp$ 侧面 $BB_1C_1C$，故 $AB \perp BE$．
在直角 $\triangle ABE$ 中，$AE^2 = AB^2 + BE^2$．
在直角 $\triangle AB_1E$ 中，已知 $EA \perp EB_1$，且 $AB \perp EB_1$，
由三垂线定理逆定理可得 $BE \perp EB_1$．
\Needspace{5\baselineskip}
因此在 $\triangle BEB_1$ 中，$\angle BEB_1 = 90^\circ$：
\nopagebreak[4]
\[
BE^2 + EB_1^2 = BB_1^2 \implies (x^2 - x + 1) + (x^2 - 5x + 7) = 2^2 = 4,
\]
\Needspace{5\baselineskip}
化简得：
\nopagebreak[4]
\[
2x^2 - 6x + 4 = 0 \implies x^2 - 3x + 2 = 0 \implies (x - 1)(x - 2) = 0.
\]
因为 $E$ 是异于 $C, C_1$ 的点，所以 $0 < x < 2$，解得 $x = 1$．
此时 $BE = \sqrt{1^2 - 1 + 1} = 1$．
因为 $BE \perp EB_1$ 且 $BE \perp AB$（由 $AB \perp$ 侧面 $BB_1C_1C$），
所以线段 $BE$ 为异面直线 $AB$ 与 $EB_1$ 的公垂线段，
因此异面直线 $AB$ 与 $EB_1$ 的距离为 $BE = 1$．

\gksubpart{(2)}

求解二面角 $A - EB_1 - A_1$ 的平面角的正切值：
在侧面 $BB_1C_1C$ 内过 $E$ 作 $EK \perp EB_1$ 交直线 $BB_1$ 于 $K$（或者直接建立空间直角坐标系）：
以 $B$ 为坐标原点：
以射线 $BA$ 的方向为 $x$ 轴正方向，单位正交基向量为 $\vec{i}$；
以射线 $BE$ 的方向为 $y$ 轴正方向，单位正交基向量为 $\vec{j}$；
以射线 $EB_1$ 的方向为 $z$ 轴正方向，单位正交基向量为 $\vec{k}$．
由 (1) 可知 $BA, BE, EB_1$ 两两垂直，$AB = \sqrt{2}, BE = 1, EB_1 = \sqrt{x^2 - 5x + 7} = \sqrt{3}$．
\Needspace{5\baselineskip}
各点坐标为：
\nopagebreak[4]
\[
B(0, 0, 0), \quad A(\sqrt{2}, 0, 0), \quad E(0, 1, 0), \quad B_1(0, 1, \sqrt{3}).
\]
又因为在三棱柱中，$\overrightarrow{AA_1} = \overrightarrow{BB_1} = (0, 1, \sqrt{3})$，
\Needspace{5\baselineskip}
所以：
\nopagebreak[4]
\[
A_1 = A + \overrightarrow{AA_1} = (\sqrt{2}, 1, \sqrt{3}).
\]
计算相关平面的法向量：
\begin{itemize}
\item 对于平面 $AEB_1$：
\end{itemize}
\nopagebreak[4]
\[
\overrightarrow{EA} = (\sqrt{2}, -1, 0), \quad \overrightarrow{EB_1} = (0, 0, \sqrt{3}).
\]
\Needspace{5\baselineskip}
设其法向量为 $\vec{n}_1 = (x_1, y_1, z_1)$，则：
\nopagebreak[4]
\[
\begin{cases}
\sqrt{2}x_1 - y_1 = 0, \\
\sqrt{3}z_1 = 0.
\end{cases}
\]
取 $x_1 = 1$，则 $y_1 = \sqrt{2}, z_1 = 0$，得 $\vec{n}_1 = (1, \sqrt{2}, 0)$；
\begin{itemize}
\item 对于平面 $A_1EB_1$：
\end{itemize}
\nopagebreak[4]
\[
\overrightarrow{B_1A_1} = (\sqrt{2}, 0, 0), \quad \overrightarrow{B_1E} = (0, 0, -\sqrt{3}).
\]
平面 $A_1EB_1$ 包含平行于 $x$ 轴与 $z$ 轴的向量，故其法向量平行于 $y$ 轴，可取 $\vec{n}_2 = (0, 1, 0)$．
\Needspace{5\baselineskip}
设二面角 $A - EB_1 - A_1$ 的平面角为 $\theta$（由图形可知 $\theta$ 为锐角），则：
\nopagebreak[4]
\[
\cos\theta = \frac{|\vec{n}_1 \cdot \vec{n}_2|}{|\vec{n}_1| |\vec{n}_2|} = \frac{|\sqrt{2}|}{\sqrt{1^2 + (\sqrt{2})^2 + 0} \times 1} = \frac{\sqrt{2}}{\sqrt{3}} = \frac{\sqrt{6}}{3}.
\]
\Needspace{5\baselineskip}
因为 $\theta$ 为锐角，所以：
\nopagebreak[4]
\[
\sin\theta = \sqrt{1 - \cos^2\theta} = \sqrt{1 - \frac{2}{3}} = \frac{\sqrt{3}}{3}.
\]
\Needspace{5\baselineskip}
因此二面角平面角的正切值为：
\nopagebreak[4]
\[
\tan\theta = \frac{\sin\theta}{\cos\theta} = \frac{\sqrt{3}/3}{\sqrt{6}/3} = \frac{1}{\sqrt{2}} = \frac{\sqrt{2}}{2}.
\]

\gkanswer{答案（1）：$1$。}

\gkanswer{答案（2）：$\dfrac{\sqrt{2}}{2}$。}

\end{answer}
\begin{answer}
\gksolhead{21}{}

\textbf{题目}：

（本小题满分 12 分）

已知椭圆 $C_1$ 的方程为 $\dfrac{x^2}{4} + y^2 = 1$，双曲线 $C_2$ 的左、右焦点分别为 $C_1$ 的左、右顶点，而 $C_2$ 的左、右顶点分别是 $C_1$ 的左、右焦点．

(1) 求双曲线 $C_2$ 的方程；

(2) 若直线 $l: y = kx + \sqrt{2}$ 与椭圆 $C_1$ 及双曲线 $C_2$ 都恒有两个不同的交点，且 $l$ 与 $C_2$ 的两个交点 $A$ 和 $B$ 满足 $\overrightarrow{OA} \cdot \overrightarrow{OB} < 6$（其中 $O$ 为原点），求 $k$ 的取值范围．

\Needspace{4\baselineskip}
\textbf{解析}

\gksubpart{(1)}

椭圆 $C_1$ 的方程为 $\dfrac{x^2}{4} + y^2 = 1$，其中 $a_1^2 = 4, b_1^2 = 1$．
半焦距 $c_1 = \sqrt{a_1^2 - b_1^2} = \sqrt{4 - 1} = \sqrt{3}$．
因此椭圆 $C_1$ 的左、右顶点为 $(\pm 2, 0)$，左、右焦点为 $(\pm\sqrt{3}, 0)$．
根据题意，双曲线 $C_2$ 的左、右焦点是 $C_1$ 的左、右顶点 $(\pm 2, 0)$，故 $c_2 = 2$；
双曲线 $C_2$ 的左、右顶点是 $C_1$ 的左、右焦点 $(\pm\sqrt{3}, 0)$，故 $a_2 = \sqrt{3}$．
所以双曲线的实半轴长 $a_2 = \sqrt{3}$，虚半轴长满足 $b_2^2 = c_2^2 - a_2^2 = 2^2 - (\sqrt{3})^2 = 4 - 3 = 1$．
\Needspace{5\baselineskip}
因为焦点在 $x$ 轴上，所以双曲线 $C_2$ 的方程为：
\nopagebreak[4]
\[
\frac{x^2}{3} - y^2 = 1.
\]

\gksubpart{(2)}

将直线 $l: y = kx + \sqrt{2}$ 分别与两曲线方程联立：
\Needspace{5\baselineskip}
首先与椭圆 $C_1$ 联立：
\nopagebreak[4]
\[
\begin{cases}
y = kx + \sqrt{2}, \\
\dfrac{x^2}{4} + y^2 = 1,
\end{cases}
\implies x^2 + 4(kx + \sqrt{2})^2 = 4 \implies (4k^2 + 1)x^2 + 8\sqrt{2}kx + 4 = 0.
\]
\Needspace{5\baselineskip}
由于直线与椭圆恒有两个不同交点，故判别式：
\nopagebreak[4]
\[
\Delta_1 = (8\sqrt{2}k)^2 - 4(4k^2 + 1) \times 4 = 128k^2 - 64k^2 - 16 = 64k^2 - 16 > 0 \iff k^2 > \frac{1}{4}.
\]
\Needspace{5\baselineskip}
其次与双曲线 $C_2$ 联立：
\nopagebreak[4]
\[
\begin{cases}
y = kx + \sqrt{2}, \\
\dfrac{x^2}{3} - y^2 = 1,
\end{cases}
\implies x^2 - 3(kx + \sqrt{2})^2 = 3 \implies (1 - 3k^2)x^2 - 6\sqrt{2}kx - 9 = 0.
\]
\Needspace{5\baselineskip}
直线与双曲线有两个不同交点，必须满足二次项系数不为 0 且判别式大于 0：
\nopagebreak[4]
\[
1 - 3k^2 \neq 0 \iff k^2 \neq \frac{1}{3},
\]
\nopagebreak[4]
\[
\Delta_2 = (-6\sqrt{2}k)^2 - 4(1 - 3k^2)(-9) = 72k^2 + 36(1 - 3k^2) = 36(1 - k^2) > 0 \iff k^2 < 1.
\]
\Needspace{5\baselineskip}
设交点 $A(x_1, y_1), B(x_2, y_2)$，由韦达定理：
\nopagebreak[4]
\[
x_1 + x_2 = \frac{6\sqrt{2}k}{1 - 3k^2}, \quad x_1 x_2 = \frac{-9}{1 - 3k^2}.
\]
\Needspace{5\baselineskip}
因为点 $A, B$ 在直线 $l$ 上，所以 $y_1 = kx_1 + \sqrt{2}, y_2 = kx_2 + \sqrt{2}$，从而：
\nopagebreak[4]
\[
y_1 y_2 = (kx_1 + \sqrt{2})(kx_2 + \sqrt{2}) = k^2 x_1 x_2 + \sqrt{2}k(x_1 + x_2) + 2
\]
\nopagebreak[4]
\[
= k^2 \left(\frac{-9}{1 - 3k^2}\right) + \sqrt{2}k \left(\frac{6\sqrt{2}k}{1 - 3k^2}\right) + 2 = \frac{3k^2}{1 - 3k^2} + 2 = \frac{2 - 3k^2}{1 - 3k^2}.
\]
\Needspace{5\baselineskip}
计算向量数量积：
\nopagebreak[4]
\[
\overrightarrow{OA} \cdot \overrightarrow{OB} = x_1 x_2 + y_1 y_2 = \frac{-9}{1 - 3k^2} + \frac{2 - 3k^2}{1 - 3k^2} = \frac{-3k^2 - 7}{1 - 3k^2}.
\]
\Needspace{5\baselineskip}
由已知条件 $\overrightarrow{OA} \cdot \overrightarrow{OB} < 6$：
\nopagebreak[4]
\[
\frac{-3k^2 - 7}{1 - 3k^2} < 6 \iff \frac{-3k^2 - 7 - 6(1 - 3k^2)}{1 - 3k^2} < 0
\]
\nopagebreak[4]
\[
\iff \frac{15k^2 - 13}{1 - 3k^2} < 0 \iff \frac{15k^2 - 13}{3k^2 - 1} > 0.
\]
\Needspace{5\baselineskip}
解此分式不等式可得：
\nopagebreak[4]
\[
k^2 < \frac{1}{3} \quad \text{\text{或}} \quad k^2 > \frac{13}{15}.
\]
\Needspace{5\baselineskip}
综合以上各条件，求 $k^2$ 的交集：
\nopagebreak[4]
\[
\begin{cases}
k^2 > \dfrac{1}{4}, \\
k^2 < 1 \text{ \text{且} } k^2 \neq \dfrac{1}{3}, \\
k^2 < \dfrac{1}{3} \text{ \text{或} } k^2 > \dfrac{13}{15}.
\end{cases}
\]
\Needspace{5\baselineskip}
解得：
\nopagebreak[4]
\[
\frac{1}{4} < k^2 < \frac{1}{3} \quad \text{\text{或}} \quad \frac{13}{15} < k^2 < 1.
\]
\Needspace{5\baselineskip}
解此不等式得 $k$ 的取值范围为：
\nopagebreak[4]
\[
\left(-1, -\frac{\sqrt{195}}{15}\right) \cup \left(-\frac{\sqrt{3}}{3}, -\frac{1}{2}\right) \cup \left(\frac{1}{2}, \frac{\sqrt{3}}{3}\right) \cup \left(\frac{\sqrt{195}}{15}, 1\right).
\]

\gkanswer{答案（1）：$\dfrac{x^2}{3} - y^2 = 1$。}

\gkanswer{答案（2）：$\left(-1, -\dfrac{\sqrt{195}}{15}\right) \cup \left(-\dfrac{\sqrt{3}}{3}, -\dfrac{1}{2}\right) \cup \left(\dfrac{1}{2}, \dfrac{\sqrt{3}}{3}\right) \cup \left(\dfrac{\sqrt{195}}{15}, 1\right)$。}

\end{answer}
\begin{answer}
\gksolhead{22}{}

\textbf{题目}：

（本小题满分 14 分）

数列 $\{a_n\}$ 满足 $a_1 = 1$ 且 $a_{n+1} = \left(1 + \dfrac{1}{n^2 + n}\right) a_n + \dfrac{1}{2^n}$ ($n \geqslant 1$)．

(1) 用数学归纳法证明：$a_n \geqslant 2$ ($n \geqslant 2$)；

(2) 已知不等式 $\ln(1 + x) < x$ 对 $x > 0$ 成立，证明：$a_n < \mathrm{e}^2$ ($n \geqslant 1$)，其中无理数 $\mathrm{e} = 2.71828\cdots$．

\Needspace{4\baselineskip}
\textbf{解析}

\gksubpart{(1)}

用数学归纳法证明 $a_n \geqslant 2$ ($n \geqslant 2$)：
① 当 $n = 2$ 时：
\Needspace{5\baselineskip}
由已知递推公式及 $a_1 = 1$：
\nopagebreak[4]
\[
a_2 = \left(1 + \frac{1}{1^2 + 1}\right) a_1 + \frac{1}{2^1} = \left(1 + \frac{1}{2}\right) \times 1 + \frac{1}{2} = \frac{3}{2} + \frac{1}{2} = 2.
\]
显然 $a_2 \geqslant 2$，结论成立；
② 假设当 $n = k$ ($k \geqslant 2, k \in \mathbf{N}^*$) 时结论成立，即 $a_k \geqslant 2$．
\Needspace{5\baselineskip}
那么当 $n = k + 1$ 时：
\nopagebreak[4]
\[
a_{k+1} = \left(1 + \frac{1}{k^2 + k}\right) a_k + \frac{1}{2^k}.
\]
因为 $k \geqslant 2$，所以 $\dfrac{1}{k^2 + k} > 0, \dfrac{1}{2^k} > 0$．
\Needspace{5\baselineskip}
由归纳假设 $a_k \geqslant 2$，得：
\nopagebreak[4]
\[
a_{k+1} \geqslant \left(1 + \frac{1}{k^2 + k}\right) \times 2 + \frac{1}{2^k} = 2 + \frac{2}{k^2 + k} + \frac{1}{2^k} > 2.
\]
即当 $n = k + 1$ 时结论也成立．
由 ①② 可知，对一切 $n \geqslant 2$，都有 $a_n \geqslant 2$．

\gksubpart{(2)}

证明 $a_n < \mathrm{e}^2$ ($n \geqslant 1$)：
当 $n = 1$ 时，$a_1 = 1 < \mathrm{e}^2$ 显然成立．
当 $n \geqslant 2$ 时，记 $p_k = 1 + \dfrac{1}{k^2 + k} > 1$ ($k \geqslant 1$)，
\Needspace{5\baselineskip}
并定义累乘积：
\nopagebreak[4]
\[
P_1 = 1, \quad P_m = \prod_{j=1}^{m-1} p_j = \prod_{j=1}^{m-1} \left(1 + \frac{1}{j^2 + j}\right) \quad (m \geqslant 2).
\]
已知递推关系为 $a_{k+1} = p_k a_k + \dfrac{1}{2^k}$．
\Needspace{5\baselineskip}
两边同除以 $P_{k+1} = p_k P_k$：
\nopagebreak[4]
\[
\frac{a_{k+1}}{P_{k+1}} = \frac{a_k}{P_k} + \frac{1}{2^k P_{k+1}}.
\]
因为对任意 $j \geqslant 1$ 均有 $p_j > 1$，所以对 $k \geqslant 1$ 恒有 $P_{k+1} > 1$，从而 $\dfrac{1}{P_{k+1}} < 1$．
\Needspace{5\baselineskip}
因此：
\nopagebreak[4]
\[
\frac{a_{k+1}}{P_{k+1}} - \frac{a_k}{P_k} < \frac{1}{2^k}.
\]
\Needspace{5\baselineskip}
对 $k = 1, 2, \dots, n-1$ 进行累加，得：
\nopagebreak[4]
\[
\frac{a_n}{P_n} - \frac{a_1}{P_1} = \sum_{k=1}^{n-1} \left(\frac{a_{k+1}}{P_{k+1}} - \frac{a_k}{P_k}\right) < \sum_{k=1}^{n-1} \frac{1}{2^k} = 1 - \frac{1}{2^{n-1}} < 1.
\]
\Needspace{5\baselineskip}
因为 $a_1 = 1, P_1 = 1$，所以 $\dfrac{a_1}{P_1} = 1$，从而：
\nopagebreak[4]
\[
\frac{a_n}{P_n} < 1 + \left(1 - \frac{1}{2^{n-1}}\right) = 2 - \frac{1}{2^{n-1}} < 2,
\]
即 $a_n < 2 P_n$．
下面估计 $P_n$ 的大小：
\Needspace{5\baselineskip}
对 $P_n$ 两边取自然对数：
\nopagebreak[4]
\[
\ln P_n = \sum_{j=1}^{n-1} \ln\left(1 + \frac{1}{j^2 + j}\right).
\]
\Needspace{5\baselineskip}
由已知条件，当 $x > 0$ 时 $\ln(1 + x) < x$，取 $x = \dfrac{1}{j^2 + j} > 0$，可得：
\nopagebreak[4]
\[
\ln\left(1 + \frac{1}{j^2 + j}\right) < \frac{1}{j^2 + j} = \frac{1}{j(j+1)} = \frac{1}{j} - \frac{1}{j+1}.
\]
\Needspace{5\baselineskip}
裂项相消求和：
\nopagebreak[4]
\[
\ln P_n < \sum_{j=1}^{n-1} \left(\frac{1}{j} - \frac{1}{j+1}\right) = 1 - \frac{1}{n} < 1.
\]
\Needspace{5\baselineskip}
两边取以 $\mathrm{e}$ 为底的指数，得：
\nopagebreak[4]
\[
P_n < \mathrm{e}^1 = \mathrm{e}.
\]
\Needspace{5\baselineskip}
结合前面的放缩不等式：
\nopagebreak[4]
\[
a_n < 2 P_n < 2\mathrm{e}.
\]
\Needspace{5\baselineskip}
因为无理数 $\mathrm{e} \approx 2.71828 > 2$，所以：
\nopagebreak[4]
\[
2\mathrm{e} < \mathrm{e} \cdot \mathrm{e} = \mathrm{e}^2.
\]
从而对一切 $n \geqslant 2$，都有 $a_n < \mathrm{e}^2$．
综上，对一切 $n \geqslant 1$，均有 $a_n < \mathrm{e}^2$．

\gkanswer{答案（1）：见解析版}

\gkanswer{答案（2）：见解析版}

\end{answer}